\section{Local normalization and its derivatives}
\label{sec:normalization-jets}
The local scalar correction is a difference between the actual sector
and its flat reference. The reference uses the same forest, tags,
cutoff radius and windows; its coarse field is the identity and its
seed sets are empty. Write $A^\circ$ for its completed precision
matrix. With the local colour trace $J_e(A)$ of
Section~\ref{sec:local-covariance}, the correction at a physical bond is
\begin{equation}\label{eq:normalization-correction}
 c_e=\frac12\{J_e(A^\circ)-J_e(A)\}.
\end{equation}
We first establish its scalar bounds and then differentiate the actual
window formula. A bound on the full determinant would not provide
these local estimates.

\subsection{Uniform and remote scalar bounds}
Let $\|A\|_\eta,\|A^\circ\|_\eta\le Q$ and let $S$ be the
uniform row and column bound for every principal inverse along
$A_s=1+s(A-1)$ and the corresponding flat interpolation, $0\le s\le1$.
The local formula is
\[
 J_e(A)=\sum_{\alpha=1}^q\int_0^1
       [(A-1)P_s]_{(e,\alpha),(e,\alpha)}\,ds.
\]
Thus $|J_e(A)|\le q(Q+1)S$, and the same holds for $A^\circ$.
By \eqref{eq:normalization-correction},
\begin{equation}\label{eq:normalization-uniform}
 |c_e|\le q(Q+1)S.
\end{equation}
Both precisions meet the hypotheses of the Gram estimate: the flat
reference has no flags, and gauge transformations preserve the flags
of the actual sector. Hence the same post-$L$ constants can be used
for the actual, flat and gauge-transformed sectors.

If the anchor of $e$ is farther than $30B$ from the seeds, choose a
single gauge making the coarse bond chords at most $Cw(x)$ on the
radius-$6B$ ball, as in \eqref{eq:mean-common-log}. Here $w$ has
upper exponent $-1/2$. For any coordinate in the window of
$(e,\alpha)$, its anchor lies within $B+2$ of that of $e$.
On these coordinates the inactive diagonal of the gauged sector
agrees with the flat reference: exactly the tree bonds are inactive.
If an entry of the selected row matrix meets such a coordinate,
either its row anchor lies within $B/4$ or both entries are zero.
In the former case that row is selected in both sectors. The gauge
transport comparison, applied to the same unit coordinate input,
gives
\[
 |M'_{pi}-M^\circ_{pi}|\le d(x),\qquad
 d(x)=C_1(B+1)w(x).
\]
The same estimate holds for the other tested coordinate. All transport
paths lie in the common $6B$ hull. Expanding the two products in the
Gram entry gives
\[
 |M'_{pi}M'_{pj}-M^\circ_{pi}M^\circ_{pj}|
 \le d(x)|M'_{pj}|+d(x)|M^\circ_{pi}|.
\]
Sum over $p$ and use the column bound $Y$ for both row matrices:
\begin{equation}\label{eq:normalization-plate-near}
 |A'_{ij}-A^\circ_{ij}|\le2Yd(x)=:\delta(x)
 \quad\hbox{on each tested principal window}.
\end{equation}
No count of all rows enters this bound.

The window contains at most $\Delta_B=m(2B+5)^4$ coordinates.
The local resolvent comparison of
Section~\ref{sec:local-covariance} now yields
\[
 |J_e(A')-J_e(A^\circ)|
 \le q\{S+(Q+1)S^2\}\Delta_B\delta(x).
\]
The complete colour trace is gauge invariant; the endpoint windows
are invariant under colour rotations. Therefore $J_e(A')=J_e(A)$,
and
\begin{equation}\label{eq:normalization-remote}
 |c_e|\le q\{S+(Q+1)S^2\}\Delta_B\delta(x)
       =:a_0(x),\qquad d(e,\text{seeds})>30B.
\end{equation}
We dropped a harmless factor $1/2$. Since $B$ has upper exponent
zero, $a_0$ has upper exponent $-1/2$.

\subsection{Matrix derivative bounds with spatial weights}
For a $C^k$ matrix-valued map $A(y)$ near $y_0$, a convenient
quantitative statement, for $f,v,\gamma\ge0$, is the existence of a nonnegative matrix $H$
such that
\begin{equation}\label{eq:matrix-jet-envelope}
 \|D^r A_{ij}(y_0)\|\le H_{ij}v^r\quad(0\le r\le k),
 \qquad \|H\|_\gamma\le f.
\end{equation}
The norm on the right is the maximum of the weighted row and column
sums. Thus it controls a common entrywise majorant of all jets, not
just the norm of each derivative separately. If $A$ and $B$ have
bounds $(f,v)$ and $(g,z)$ with $g,z\ge0$, Leibniz's rule gives the envelope $HK$
for their product with bounds $(fg,v+z)$: the binomial sum is
$(v+z)^r$, and the triangle inequality for distance proves
$\|HK\|_\gamma\le\|H\|_\gamma\|K\|_\gamma$.
Addition, transpose and restriction to a principal set preserve the
corresponding bounds. These operations introduce no volume factor.

Suppose $A(y_0)$ is invertible and its inverse has weighted price
$g$. Set $F=A(y_0)^{-1}$ and $V(y)=1-FA(y)$. Then $V(y_0)=0$,
and, wherever $A(y)$ is invertible,
\begin{equation}\label{eq:inverse-finite-jet}
 A(y)^{-1}=\sum_{r=0}^k V(y)^rF+V(y)^{k+1}A(y)^{-1}.
\end{equation}
Multiplying the finite geometric identity on the right by $A(y)^{-1}$
proves this without commuting any factors. Invertibility holds on a
neighbourhood by continuity of the determinant. The inverse is $C^k$
there by the adjugate formula. Every derivative of order at most $k$
of $V^{k+1}$ vanishes at $y_0$: in each Leibniz term at least one of
the $k+1$ factors is undifferentiated and equals zero. The remainder
in \eqref{eq:inverse-finite-jet} consequently has zero jets through
order $k$.

The envelope for $V$ has bounds $(1+fg,v)$. Applying the product
rules to the finite sum proves an inverse envelope with bounds
\begin{equation}\label{eq:inverse-jet-price}
 T_k(f,g)=\sum_{r=0}^k(1+fg)^r g,\qquad kv.
\end{equation}
Only the inverse at the base point was priced. A uniform inverse
bound on an entire neighbourhood was not assumed.

Apply this argument separately in every prescribed principal window.
Each restriction of $A(y)$ retains the envelope $(f,v)$, and each
base inverse has weighted price $S$. Its $(i,j)$ entry has majorant
$T_k(f,S)e^{-\gamma d(i,j)}$. Padding it by zero outside the window
gives a common envelope for $P(y)$, because window membership is fixed.
Both its row and column sums are bounded by
\begin{equation}\label{eq:pad-jet-price}
 F_B=\Delta_B T_k(f,S),\qquad\hbox{with derivative scale }kv.
\end{equation}
For the row bound, even though different columns use different
windows, any nonzero entry has $d(i,j)\le B+2$, and there are at
most $\Delta_B$ such $j$. The same count proves the column bound.

The error $A(y)P(y)-1$ vanishes, as a germ, whenever $d(i,j)\le B$.
Indeed each such coordinate belongs to the window, and its principal
inverse solves the corresponding row equation for all nearby $y$.
Erasing these zero jets and decreasing the spatial weight from
$\gamma$ to $\gamma/2$ gives an envelope with price
\begin{equation}\label{eq:inverse-error-jets}
 e^{-\gamma B/2}(fF_B+1),\qquad
 \hbox{derivative scale }v+kv.
\end{equation}
This is the differentiated version of the boundary error. Multiplying
by the full inverse envelope also prices $P-A^{-1}$; taking the
symmetric part and then multiplying by $A$ prices the covariance
defect $A(P+P^*)/2-1$. For fixed $k$ all these prices are a polynomial
in $B$ times $e^{-\gamma B/2}$, hence have every prescribed negative
upper exponent. The derivative scales are fixed multiples of $v$.

\subsection{The actual sector matrix supplies the premise}
Freeze the seeds, tree, tags, windows and active sets while varying
the coarse field $U(y)$. Suppose every link matrix has jets
\[
 \|D^r U_e(y_0)\|_{\rm F}\le R_G v^r,
 \qquad R_G=\sqrt{\dim(\text{matrix representation})},
 \quad0\le r\le k.
\]
The same bound holds for inverses by taking the adjoint, for signed
pullbacks, and for replicas. Constant gauge transformations preserve
it by unitary invariance of the Frobenius norm.

A product along a walk of length $n$ has amplitude $R_G$ and
derivative scale $nR_Gv$. To verify the absence of an exponential
amplitude, split its first factor as $U(y_0)+(U(y)-U(y_0))$.
Multiplication by $U(y_0)$ is an isometry; in every positive-order
term from the difference, its bound $R_Gv^r$ is at most
$(R_Gv)^r$. The binomial formula therefore adds $R_Gv$ to the
derivative scale while preserving the amplitude of the remaining
product. Induction proves the claim, including inverse letters and
repeated physical links.

For the selected sector row matrix, a single coordinate input is
transported from its physical bond to the row anchor. The coordinate
path has length at most $4B$. A reversed signed bond adds its replica
adjoint, whose derivative scale is $2v$. The two sides of the
transported adjoint add twice the path scale. Thus the transported
input has fixed amplitude and derivative scale
\begin{equation}\label{eq:normalization-kernel-growth}
 z=8BR_Gv+2v.
\end{equation}
The remaining cut flat row is a fixed linear map. The weighted impulse
estimate from Section~\ref{sec:sector-gram}, before periodization,
bounds its operator norm by $C e^{-\eta d(p,i)}$. The extra one-step
shift from a reversed bond is absorbed into this fixed $C$. Hence
\[
 \|D^rM_{pi}(y_0)\|\le C e^{-\eta d(p,i)}z^r.
\]
Inactive columns, unselected rows and entries outside the cutoff are
identically zero throughout the variation; their indicators are not
differentiated.

For $0<\gamma<\eta$, sum the displayed envelope using the bounded
number of row and column coordinates per coarse site. The shell sum
$\sum_{n\ge0}(2n+1)^4e^{-(\eta-\gamma)n}$ is finite. It follows
that $M$ has an envelope $(C_\gamma,z)$, with $C_\gamma$ independent
of the volume and of $B$. The identity
$A=M^*M+\Pi_{\rm inactive}$ therefore supplies
\eqref{eq:matrix-jet-envelope} with
\[
 f=C_\gamma^2+1,\qquad v_A=2z.
\]
The inactive projection is constant. If the incoming link scale $v$
has upper exponent $-1/2$, so do $z$ and $v_A$.

\subsection{Differentiating the local logarithmic trace}
Choose $\gamma$ also below the retained inverse width. The glide
$A_s(y)=1+s(A(y)-1)$ has envelope $(f+2,v_A)$ uniformly for
$0\le s\le1$. Its principal inverses at $y_0$ satisfy the same
uniform price $S$ as before. Multiplication by $A(y)-1$ and
\eqref{eq:pad-jet-price} consequently give, for every local diagonal
entry, the amplitude
\[
 (f+1)\Delta_BT_k(f+2,S)
\]
and derivative scale $(k+1)v_A$, uniformly in $s$.

To justify differentiation under the integral, note that the
principal determinants are nonzero at every $(y_0,s)$ on the compact
interval. There are finitely many windows. Joint continuity gives
one neighbourhood of $y_0$ on which all these determinants remain
nonzero for all $s\in[0,1]$. The adjugate formula makes the integrand
jointly $C^k$ there. On a smaller closed parameter ball times the
interval, all derivatives are bounded by compactness. Difference
quotients and the fundamental theorem of calculus, followed by
dominated convergence, now permit differentiation under the integral,
successively through order $k$. For this argument the variation space
is finite dimensional, as it is for each finite-volume field variation.
The local bounds at $y_0$ are unaffected by the auxiliary compactness
bound. Summing the $q$ colour entries proves
\begin{equation}\label{eq:normalization-log-jets}
 \|D^rJ_e(A(y_0))\|
 \le q(f+1)\Delta_BT_k(f+2,S)\{(k+1)v_A\}^r,
 \quad0\le r\le k.
\end{equation}

The flat-reference term in \eqref{eq:normalization-correction} is
constant in $y$. Combining its scalar bound with
\eqref{eq:normalization-log-jets}, and enlarging the harmless factor
$1/2$, gives an envelope for $c_e(y)$ with amplitude $b(x)$ of upper
exponent zero and derivative scale $z_1(x)$ of upper exponent $-1/2$.
It follows that all its jets are bounded by
$b(x)(1+z_1(x))^k$, again of upper exponent zero. At a remote bond
we can do better: its value is at most $a_0(x)$ by
\eqref{eq:normalization-remote}, while for positive orders, once
$0\le z_1\le1$,
\[
 \|D^r c_e(y_0)\|\le b(x)z_1(x)^r\le b(x)z_1(x).
\]
Thus the common bound $a_0+bz_1$, of upper exponent $-1/2$, controls
all jets through order $k$ at remote bonds. This obtains small
normalization jets without differentiating a chosen local gauge.
The needed link variation and its coherent Gaussian compensation
remain distinct inputs to the eventual packet variation.
